\documentclass{amsart}
% For dealing with references we use the comment environment
\usepackage{verbatim}
\newenvironment{reference}{\comment}{\endcomment}
%\newenvironment{reference}{}{}
\newenvironment{slogan}{\comment}{\endcomment}
\newenvironment{history}{\comment}{\endcomment}

% For commutative diagrams we use Xy-pic
\usepackage[all]{xy}

% We use 2cell for 2-commutative diagrams.
\xyoption{2cell}
\UseAllTwocells

% We use multicol for the list of chapters between chapters
\usepackage{multicol}

% This is generally recommended for better output
\usepackage{lmodern}
\usepackage[T1]{fontenc}

% For cross-file-references

% Package for hypertext links:
\usepackage{hyperref}

% For any local file, say "hello.tex" you want to link to please

% Theorem environments.
%
\theoremstyle{plain}
\newtheorem{theorem}[subsection]{Theorem}
\newtheorem{proposition}[subsection]{Proposition}
\newtheorem{lemma}[subsection]{Lemma}

\theoremstyle{definition}
\newtheorem{definition}[subsection]{Definition}
\newtheorem{example}[subsection]{Example}
\newtheorem{exercise}[subsection]{Exercise}
\newtheorem{situation}[subsection]{Situation}

\theoremstyle{remark}
\newtheorem{remark}[subsection]{Remark}
\newtheorem{remarks}[subsection]{Remarks}

\numberwithin{equation}{subsection}

% Macros
%
\def\lim{\mathop{\mathrm{lim}}\nolimits}
\def\colim{\mathop{\mathrm{colim}}\nolimits}
\def\Spec{\mathop{\mathrm{Spec}}}
\def\Hom{\mathop{\mathrm{Hom}}\nolimits}
\def\Ext{\mathop{\mathrm{Ext}}\nolimits}
\def\SheafHom{\mathop{\mathcal{H}\!\mathit{om}}\nolimits}
\def\SheafExt{\mathop{\mathcal{E}\!\mathit{xt}}\nolimits}
\def\Sch{\mathit{Sch}}
\def\Mor{\mathop{\mathrm{Mor}}\nolimits}
\def\Ob{\mathop{\mathrm{Ob}}\nolimits}
\def\Sh{\mathop{\mathit{Sh}}\nolimits}
\def\NL{\mathop{N\!L}\nolimits}
\def\CH{\mathop{\mathrm{CH}}\nolimits}
\def\proetale{{pro\text{-}\acute{e}tale}}
\def\etale{{\acute{e}tale}}
\def\QCoh{\mathit{QCoh}}
\def\Ker{\mathop{\mathrm{Ker}}}
\def\Im{\mathop{\mathrm{Im}}}
\def\Coker{\mathop{\mathrm{Coker}}}
\def\Coim{\mathop{\mathrm{Coim}}}

% Boxtimes
%
\DeclareMathSymbol{\boxtimes}{\mathbin}{AMSa}{"02}

%
% Macros for moduli stacks/spaces
%
\def\QCohstack{\mathcal{QC}\!\mathit{oh}}
\def\Cohstack{\mathcal{C}\!\mathit{oh}}
\def\Spacesstack{\mathcal{S}\!\mathit{paces}}
\def\Quotfunctor{\mathrm{Quot}}
\def\Hilbfunctor{\mathrm{Hilb}}
\def\Curvesstack{\mathcal{C}\!\mathit{urves}}
\def\Polarizedstack{\mathcal{P}\!\mathit{olarized}}
\def\Complexesstack{\mathcal{C}\!\mathit{omplexes}}
% \Pic is the operator that assigns to X its picard group, usage \Pic(X)
% \Picardstack_{X/B} denotes the Picard stack of X over B
% \Picardfunctor_{X/B} denotes the Picard functor of X over B
\def\Pic{\mathop{\mathrm{Pic}}\nolimits}
\def\Picardstack{\mathcal{P}\!\mathit{ic}}
\def\Picardfunctor{\mathrm{Pic}}
\def\Deformationcategory{\mathcal{D}\!\mathit{ef}}

\setlength{\emergencystretch}{3em}
\begin{document}
\title[Stacks Project, Tag 07CP]{730 --- Mod-square algebra lifts preserving smooth loci under stable annihilators over Noetherian rings}
\maketitle
\noindent Copyright (C) 2005--2025 Johan de Jong. Stacks Project authors. Source: \href{https://stacks.math.columbia.edu/tag/07CP}{Tag 07CP}.

\noindent Editorial difficulty note (not author text). Difficulty H3: The author gives an explicit presentation with auxiliary variables and compatibility relations, proves generic smoothness, and identifies local models using annihilator stability. The lengthy gluing of presentation data and two independent smooth-locus conclusions is specialized desingularization technology..

\noindent Editorial provenance note (not author text). History: excerpt from revision \texttt{a0444\allowbreak 6e57e\allowbreak c1fbc\allowbreak 252a8\allowbreak 71afc\allowbreak ec775\allowbreak 2fb28\allowbreak 07b14}, smoothing.tex, lines 1420--1651. Reference presentation was adapted; original-author context and core support blocks were retained or added. The mathematical wording is unchanged. Licensed under GNU FDL 1.2 or later; no invariant sections or cover texts. See ../licenses/COPYING. Context blocks reproduce author definitions and supporting results; they are not additional counted problems.

\begin{definition}
\label{algebra-definition-smooth}
A ring map $R \to S$ is {\it smooth} if it is of finite presentation
and the naive cotangent complex $\NL_{S/R}$ is quasi-isomorphic to a
finite projective $S$-module placed in degree $0$: this means
that $H_1(\NL_{S/R}) = 0$ and that $\Omega_{S/R}$ is a finite projective
$S$-module.
\end{definition}

\begin{definition}
\label{algebra-definition-smooth-at-prime}
Let $R \to S$ be a ring map.
Let $\mathfrak q$ be a prime of $S$.
We say $R \to S$ is {\it smooth at $\mathfrak q$} if there
exists a $g \in S$, $g \not \in \mathfrak q$ such
that $R \to S_g$ is smooth.
\end{definition}

\begin{lemma}
\label{lemma-find-strictly-standard}
Let $R$ be a ring. Let $A = R[x_1, \ldots, x_n]/(f_1, \ldots, f_m)$.
Let $\mathfrak q \subset A$ be a prime ideal. Assume $R \to A$ is smooth
at $\mathfrak q$. Then there exists an $a \in A$, $a \not \in \mathfrak q$,
an integer $c$, $0 \leq c \leq \min(n, m)$, subsets
$U \subset \{1, \ldots, n\}$, $V \subset \{1, \ldots, m\}$
of cardinality $c$ such that
$$
a = a' \det(\partial f_j/\partial x_i)_{j \in V, i \in U}
$$
for some $a' \in A$ and
$$
a f_\ell \in (f_j, j \in V) + (f_1, \ldots, f_m)^2
$$
for all $\ell \in \{1, \ldots, m\}$.
\end{lemma}

\begin{proof}
Set $I = (f_1, \ldots, f_m)$ so that the naive cotangent
complex of $A$ over $R$ is homotopy equivalent to
$I/I^2 \to \bigoplus A\text{d}x_i$, see
Algebra, Lemma \href{https://stacks.math.columbia.edu/tag/00S1}{lemma NL homotopy (Tag 00S1)}.
We will use the formation of the naive cotangent complex commutes with
localization, see Algebra, Section \href{https://stacks.math.columbia.edu/tag/00S0}{section netherlander (Tag 00S0)},
especially Algebra, Lemma \href{https://stacks.math.columbia.edu/tag/00S7}{lemma localize NL (Tag 00S7)}.
By Algebra, Definitions \ref{algebra-definition-smooth} and
\ref{algebra-definition-smooth-at-prime}
we see that $(I/I^2)_a \to \bigoplus A_a\text{d}x_i$
is a split injection for some $a \in A$, $a \not \in \mathfrak q$.
After renumbering $x_1, \ldots, x_n$ and $f_1, \ldots, f_m$ we may
assume that $f_1, \ldots, f_c$ form a basis for
the vector space $I/I^2 \otimes_A \kappa(\mathfrak q)$ and that
$\text{d}x_{c + 1}, \ldots, \text{d}x_n$ map to a basis of
$\Omega_{A/R} \otimes_A \kappa(\mathfrak q)$. Hence after replacing $a$
by $aa'$ for some $a' \in A$, $a' \not \in \mathfrak q$ we may assume
$f_1, \ldots, f_c$ form a basis for $(I/I^2)_a$ and that
$\text{d}x_{c + 1}, \ldots, \text{d}x_n$ map to a basis of
$(\Omega_{A/R})_a$. In this situation $a^N$ for some large integer
$N$ satisfies the conditions of the lemma (with $U = V = \{1, \ldots, c\}$).
\end{proof}


\begin{lemma}
\label{lemma-lifting}
Let $R$ be a Noetherian ring. Let $\Lambda$ be an $R$-algebra.
Let $\pi \in R$ and assume that $\text{Ann}_R(\pi) = \text{Ann}_R(\pi^2)$ and
$\text{Ann}_\Lambda(\pi) = \text{Ann}_\Lambda(\pi^2)$.
Suppose we have $R$-algebra maps
$R/\pi^2R \to \bar C \to \Lambda/\pi^2\Lambda$
with $\bar C$ of finite presentation.
Then there exists an $R$-algebra homomorphism
$D \to \Lambda$ and a commutative diagram
$$
\xymatrix{
R/\pi^2R \ar[r] \ar[d] &
\bar C \ar[r] \ar[d] &
\Lambda/\pi^2\Lambda \ar[d] \\
R/\pi R \ar[r] &
D/\pi D \ar[r] &
\Lambda/\pi \Lambda
}
$$
with the following properties
\begin{enumerate}
\item[(a)] $D$ is of finite presentation,
\item[(b)] $R \to D$ is smooth at any prime $\mathfrak q$ with
$\pi \not \in \mathfrak q$,
\item[(c)] $R \to D$ is smooth at any prime $\mathfrak q$ with
$\pi \in \mathfrak q$ lying over a prime of $\bar C$ where
$R/\pi^2 R \to \bar C$ is smooth, and
\item[(d)] $\bar C/\pi \bar C \to D/\pi D$ is smooth at any prime
lying over a prime of $\bar C$ where $R/\pi^2R \to \bar C$ is smooth.
\end{enumerate}
\end{lemma}

\begin{proof}
We choose a presentation
$$
\bar C = R[x_1, \ldots, x_n]/(f_1, \ldots, f_m)
$$
We also denote $I = (f_1, \ldots, f_m)$ and $\bar I$ the image of
$I$ in $R/\pi^2R[x_1, \ldots, x_n]$. Since $R$ is Noetherian, so is
$\bar C$. Hence the smooth locus of $R/\pi^2 R \to \bar C$
is quasi-compact, see
Topology, Lemma \href{https://stacks.math.columbia.edu/tag/0052}{lemma Noetherian (Tag 0052)}.
Applying
Lemma \ref{lemma-find-strictly-standard}
we may choose a finite list of elements
$a_1, \ldots, a_r \in R[x_1, \ldots, x_n]$ such that
\begin{enumerate}
\item the union of the open subspaces
$\Spec(\bar C_{a_k}) \subset \Spec(\bar C)$
cover the smooth locus of $R/\pi^2 R \to \bar C$, and
\item for each $k = 1, \ldots, r$ there exists a finite subset
$E_k \subset \{1, \ldots, m\}$ such that
$(\bar I/\bar I^2)_{a_k}$ is freely generated by the classes of
$f_j$, $j \in E_k$.
\end{enumerate}
Set $I_k = (f_j, j \in E_k) \subset I$ and denote $\bar I_k$ the
image of $I_k$ in $R/\pi^2R[x_1, \ldots, x_n]$.
By (2) and Nakayama's lemma we see that $(\bar I/\bar I_k)_{a_k}$
is annihilated by $1 + b'_k$ for some $b'_k \in \bar I_{a_k}$.
Suppose $b'_k$ is the image of $b_k/(a_k)^N$ for some $b_k \in I$
and some integer $N$. After replacing $a_k$ by $a_kb_k$ we get
\begin{enumerate}
\item[(3)] $(\bar I_k)_{a_k} = (\bar I)_{a_k}$.
\end{enumerate}
Thus, after possibly replacing $a_k$ by a high power, we may write
\begin{enumerate}
\item[(4)]
$a_k f_\ell = \sum\nolimits_{j \in E_k} h_{k, \ell}^jf_j + \pi^2 g_{k, \ell}$
\end{enumerate}
for any $\ell \in \{1, \ldots, m\}$ and some
$h_{i, \ell}^j, g_{i, \ell} \in R[x_1, \ldots, x_n]$.
If $\ell \in E_k$ we choose $h_{k, \ell}^j = a_k\delta_{\ell, j}$
(Kronecker delta) and $g_{k, \ell} = 0$. Set
$$
D = R[x_1, \ldots, x_n, z_1, \ldots, z_m]/
(f_j - \pi z_j, p_{k, \ell}).
$$
Here $j \in \{1, \ldots, m\}$, $k \in \{1, \ldots, r\}$,
$\ell \in \{1, \ldots, m\}$, and
$$
p_{k, \ell} = a_k z_\ell - \sum\nolimits_{j \in E_k} h_{k, \ell}^j z_j
- \pi g_{k, \ell}.
$$
Note that for $\ell \in E_k$ we have $p_{k, \ell} = 0$ by our choices above.

\medskip\noindent
The map $R \to D$ is the given one.
Say $\bar C \to \Lambda/\pi^2\Lambda$ maps $x_i$
to the class of $\lambda_i$ modulo $\pi^2$. For an element
$f \in R[x_1, \ldots, x_n]$ we denote $f(\lambda) \in \Lambda$
the result of substituting $\lambda_i$ for $x_i$. Then we know that
$f_j(\lambda) = \pi^2 \mu_j$ for some $\mu_j \in \Lambda$.
Define $D \to \Lambda$ by the rules $x_i \mapsto \lambda_i$ and
$z_j \mapsto \pi\mu_j$. This is well defined because
\begin{align*}
p_{k, \ell} & \mapsto
a_k(\lambda) \pi \mu_\ell -
\sum\nolimits_{j \in E_k} h_{k, \ell}^j(\lambda) \pi \mu_j
- \pi g_{k, \ell}(\lambda) \\
& =
\pi\left(a_k(\lambda) \mu_\ell -
\sum\nolimits_{j \in E_k} h_{k, \ell}^j(\lambda) \mu_j
- g_{k, \ell}(\lambda)\right)
\end{align*}
Substituting $x_i = \lambda_i$ in (4) above we see that the expression
inside the brackets is annihilated by $\pi^2$, hence it is annihilated
by $\pi$ as we have assumed
$\text{Ann}_\Lambda(\pi) = \text{Ann}_\Lambda(\pi^2)$.
The map $\bar C \to D/\pi D$ is determined by $x_i \mapsto x_i$
(clearly well defined). Thus we are done if we can prove (b), (c), and (d).

\medskip\noindent
Using (4) we obtain the following key equality
\begin{align*}
\pi p_{k, \ell} & =
\pi a_k z_\ell - \sum\nolimits_{j \in E_k} \pi h_{k, \ell}^jz_j
- \pi^2 g_{k, \ell} \\
& =
- a_k (f_\ell - \pi z_\ell) + a_k f_\ell +
\sum\nolimits_{j \in E_k} h_{k, \ell}^j (f_j - \pi z_j) -
\sum\nolimits_{j \in E_k} h_{k, \ell}^j f_j - \pi^2 g_{k, \ell} \\
& =
-a_k(f_\ell - \pi z_\ell) +
\sum\nolimits_{j \in E_k} h_{k, \ell}^j(f_j - \pi z_j)
\end{align*}
The end result is an element of the ideal generated by $f_j - \pi z_j$.
In particular, we see that $D[1/\pi]$ is isomorphic to
$R[1/\pi][x_1, \ldots, x_n, z_1, \ldots, z_m]/(f_j - \pi z_j)$
which is isomorphic to $R[1/\pi][x_1, \ldots, x_n]$ hence smooth
over $R$. This proves (b).

\medskip\noindent
For fixed $k \in \{1, \ldots, r\}$ consider the ring
$$
D_k = R[x_1, \ldots, x_n, z_1, \ldots, z_m]/
(f_j - \pi z_j, j \in E_k, p_{k, \ell})
$$
The number of equations is $m = |E_k| + (m - |E_k|)$ as $p_{k, \ell}$
is zero if $\ell \in E_k$. Also, note that
\begin{align*}
(D_k/\pi D_k)_{a_k}
& =
R/\pi R[x_1, \ldots, x_n, 1/a_k, z_1, \ldots, z_m]/
(f_j, j \in E_k, p_{k, \ell}) \\
& =
(\bar C/\pi \bar C)_{a_k}[z_1, \ldots, z_m]/
(a_kz_\ell - \sum\nolimits_{j \in E_k} h_{k, \ell}^j z_j) \\
& \cong
(\bar C/\pi \bar C)_{a_k}[z_j, j \in E_k]
\end{align*}
In particular $(D_k/\pi D_k)_{a_k}$ is smooth over $(\bar C/\pi \bar C)_{a_k}$.
By our choice of $a_k$ we have that $(\bar C/\pi \bar C)_{a_k}$ is smooth
over $R/\pi R$ of relative dimension $n - |E_k|$, see (2). Hence for a prime
$\mathfrak q_k \subset D_k$ containing $\pi$ and lying over
$\Spec(\bar C_{a_k})$ the fibre ring of $R \to D_k$
is smooth at $\mathfrak q_k$ of dimension $n$. Thus $R \to D_k$ is syntomic
at $\mathfrak q_k$ by our count of the number of equations above, see
Algebra, Lemma \href{https://stacks.math.columbia.edu/tag/00ST}{lemma localize relative complete intersection (Tag 00ST)}.
Hence $R \to D_k$ is smooth at $\mathfrak q_k$, see
Algebra, Lemma \href{https://stacks.math.columbia.edu/tag/00TF}{lemma flat fibre smooth (Tag 00TF)}.

\medskip\noindent
To finish the proof, let $\mathfrak q \subset D$ be a prime
containing $\pi$ lying over a prime where $R/\pi^2 R \to \bar C$
is smooth. Then $a_k \not \in \mathfrak q$ for some $k$ by (1).
We will show that the surjection $D_k \to D$ induces
an isomorphism on local rings at $\mathfrak q$. Since we know that
the ring maps $\bar C/\pi \bar C \to D_k/\pi D_k$ and
$R \to D_k$ are smooth at the corresponding prime $\mathfrak q_k$
by the preceding paragraph this will prove (c) and (d) and thus
finish the proof.

\medskip\noindent
First, note that for any $\ell$ the equation
$\pi p_{k, \ell} = -a_k(f_\ell - \pi z_\ell) +
\sum_{j \in E_k} h_{k, \ell}^j (f_j - \pi z_j)$ proved above shows that
$f_\ell - \pi z_\ell$ maps to zero in $(D_k)_{a_k}$ and in particular
in $(D_k)_{\mathfrak q_k}$.
The relations (4) imply that $a_k f_\ell =
\sum_{j \in E_k} h_{k, \ell}^j f_j$ in $I/I^2$.
Since $(\bar I_k/\bar I_k^2)_{a_k}$ is free on $f_j$, $j \in E_k$
we see that
$$
a_{k'} h_{k, \ell}^j -
\sum\nolimits_{j' \in E_{k'}} h_{k', \ell}^{j'} h_{k, j'}^j
$$
is zero in $\bar C_{a_k}$ for every $k, k', \ell$ and $j \in E_k$.
Hence we can find a large integer $N$ such that
$$
a_k^N\left(
a_{k'} h_{k, \ell}^j -
\sum\nolimits_{j' \in E_{k'}} h_{k', \ell}^{j'} h_{k, j'}^j
\right)
$$
is in $I_k + \pi^2R[x_1, \ldots, x_n]$. Computing modulo $\pi$ we have
\begin{align*}
&
a_kp_{k', \ell} - a_{k'}p_{k, \ell} + \sum h_{k', \ell}^{j'} p_{k, j'}
\\
&
=
- a_k \sum h_{k', \ell}^{j'} z_{j'}
+ a_{k'} \sum h_{k, \ell}^j z_j
+ \sum h_{k', \ell}^{j'} a_k z_{j'}
- \sum \sum h_{k', \ell}^{j'} h_{k, j'}^j z_j \\
&
=
\sum \left(
a_{k'} h_{k, \ell}^j
- \sum h_{k', \ell}^{j'} h_{k, j'}^j
\right) z_j
\end{align*}
with Einstein summation convention. Combining with the above we see
$a_k^{N + 1} p_{k', \ell}$ is contained in the ideal generated
by $I_k$ and $\pi$ in $R[x_1, \ldots, x_n, z_1, \ldots, z_m]$.
Thus $p_{k', \ell}$ maps into $\pi (D_k)_{a_k}$. On the other hand,
the equation
$$
\pi p_{k', \ell} =
-a_{k'} (f_\ell - \pi z_\ell) +
\sum\nolimits_{j' \in E_{k'}} h_{k', \ell}^{j'}(f_{j'} - \pi z_{j'})
$$
shows that $\pi p_{k', \ell}$ is zero in $(D_k)_{a_k}$.
Since we have assumed that $\text{Ann}_R(\pi) = \text{Ann}_R(\pi^2)$
and since $(D_k)_{\mathfrak q_k}$ is smooth hence flat over $R$
we see that
$\text{Ann}_{(D_k)_{\mathfrak q_k}}(\pi) =
\text{Ann}_{(D_k)_{\mathfrak q_k}}(\pi^2)$.
We conclude that $p_{k', \ell}$ maps to zero as well, hence
$D_{\mathfrak q} = (D_k)_{\mathfrak q_k}$ and we win.
\end{proof}
\end{document}
